\documentclass[12pt,a4paper]{article}

\usepackage[utf8]{inputenc}
\usepackage[T1]{fontenc}
\usepackage{amsmath,amssymb,amsthm}
\usepackage{mathtools}
\usepackage{hyperref}
\usepackage[margin=1in]{geometry}
\usepackage{enumitem}
\usepackage{xcolor}
\usepackage{tcolorbox}

\tcbuselibrary{breakable,skins}

\newtcolorbox{mathblock}{
  colback=blue!3,
  colframe=blue!40,
  leftrule=4pt,
  rightrule=0pt,
  toprule=0pt,
  bottomrule=0pt,
  boxsep=5pt,
  breakable
}

\newtcolorbox{definitionbox}{
  colback=blue!8,
  colframe=blue!60,
  leftrule=5pt,
  rightrule=0pt,
  toprule=0pt,
  bottomrule=0pt,
  boxsep=5pt,
  breakable
}

\newtcolorbox{keypoint}{
  colback=yellow!10,
  colframe=orange!70,
  leftrule=5pt,
  rightrule=0pt,
  toprule=0pt,
  bottomrule=0pt,
  boxsep=5pt,
  breakable
}

\newcommand{\HH}{\mathbb{H}}
\newcommand{\RR}{\mathbb{R}}
\newcommand{\ZZ}{\mathbb{Z}}
\newcommand{\ii}{\mathbf{i}}
\newcommand{\jj}{\mathbf{j}}
\newcommand{\kk}{\mathbf{k}}

\hypersetup{
  colorlinks=true,
  linkcolor=blue!60!black,
  urlcolor=blue!60!black
}

\title{Quaternion-Lifted Geometric Riemann Hypothesis\\[6pt]
\large Complete Reformulation \& Analytic Bridging}
\author{}
\date{April 2026}

\begin{document}

\maketitle

\begin{center}
\textit{Merged \& rigorously unified document combining the full geometric quaternion\\
reformulation with the updated analytic bridging via the new Tao--Yang--Guth\\
zero-density breakthrough}
\end{center}

\begin{tcolorbox}[colback=red!5,colframe=red!70,leftrule=6pt,rightrule=2pt,toprule=2pt,bottomrule=2pt,boxsep=5pt]
\textbf{\textcolor{red!70!black}{\large $\triangle$ Disclaimer.}} This document is a \emph{research exposition} presenting a heuristic geometric framework for studying the Riemann Hypothesis. It does \textbf{not} constitute a proof or claimed proof of RH. Several key steps are explicitly identified as conjectural, heuristic, or open problems. The framework is offered as a potential research direction, not as a finished argument.
\end{tcolorbox}

\tableofcontents

\newpage

%----------------------------------------------------------------------
\section{Introduction \& Motivation}
\label{sec:intro}

The original geometric intuition for the Riemann Hypothesis begins with a continuous trigonometric detector of integer divisors. We reformulate this detector rigorously in the quaternion algebra~$\HH$, define a natural quaternion-valued energy functional, and lift it to the critical strip via the imaginary parts of potential zeros. The resulting conjecture states that perfect phase alignment (vanishing energy) occurs if and only if the zero lies on the critical line $\Re(s) = 1/2$.

All definitions below are precise and self-contained. We work throughout in the quaternion algebra
\[
\HH = \{ a + b\ii + c\jj + d\kk \mid a,b,c,d\in\RR \}
\]
with the standard multiplication rules $\ii^2 = \jj^2 = \kk^2 = -1$, $\ii\jj=\kk$, etc.

%----------------------------------------------------------------------
\section{The Classical Cosine Divisor Detector}
\label{sec:divisor}

For a fixed positive integer~$p$, the function

\begin{mathblock}
\[
f(x) = \cos\!\left(2\pi \frac{p}{x}\right) + \cos(2\pi x) - 2, \qquad x > 0
\]
\end{mathblock}

\noindent satisfies $f(x) = 0$ if and only if both arguments are integers, i.e., precisely when $x$ divides~$p$. This is because $\cos(2\pi \theta) = 1$ if and only if $\theta \in \ZZ$.

\begin{definitionbox}
\textbf{Divisor Detector Property (rigorous).}
Let $D_p = \{ d \in \ZZ^+ : d \mid p \}$. Then
\[
f(x) = 0 \iff x \in D_p \cup \{ p/d : d \in D_p \}.
\]
The function is smooth on $\RR^+$ and vanishes exactly on the hyperbolic branches $xy = p$ at integer lattice points.
\end{definitionbox}

%----------------------------------------------------------------------
\section{Quaternion Lifting of the Detector}
\label{sec:quaternion-lift}

We separate the two cosine terms into orthogonal imaginary directions of~$\HH$:

\begin{mathblock}
\[
\phi_p(x) := \cos(2\pi x)\,\ii + \cos\!\left(2\pi \frac{p}{x}\right)\,\jj \in \operatorname{Im}\HH.
\]
\end{mathblock}

\noindent The quaternion norm recovers the classical detector:

\begin{mathblock}
\[
\|\phi_p(x) - (\ii + \jj)\|^2 = 2f(x) + 4.
\]
\end{mathblock}

\noindent Thus perfect alignment $\phi_p(x) = \ii + \jj$ occurs precisely at the divisors of~$p$.

%----------------------------------------------------------------------
\section{The Quaternion Energy Functional}
\label{sec:energy}

Define the squared Euclidean distance in the imaginary quaternion plane:

\begin{mathblock}
\[
E_p(x) := \|\phi_p(x) - (\ii + \jj)\|^2 = 2\bigl[1 - \cos(2\pi x)\bigr]^2 + 2\bigl[1 - \cos(2\pi p/x)\bigr]^2.
\]
\end{mathblock}

\noindent Equivalently (up to the constant factor~2):

\begin{mathblock}
\[
E_p(x) = 2 \bigl[ (1 - \cos(2\pi x))^2 + (1 - \cos(2\pi p/x))^2 \bigr].
\]
\end{mathblock}

\noindent This is the natural quaternion-lifted energy functional. It vanishes exactly on the divisor set~$D_p$.

\begin{definitionbox}
\textbf{Remark on the role of quaternions.} The quaternion algebra~$\HH$ is used here purely for geometric convenience: the two cosine terms are placed on orthogonal imaginary axes $\ii$ and~$\jj$, and the energy functional is the squared Euclidean distance in $\operatorname{Im}\HH \cong \RR^3$. No non-commutative property of~$\HH$ is exploited --- the $\kk$-component is always zero, and the entire construction could equivalently be formulated in~$\RR^2$. The quaternion language is retained for consistency with the original formulation and because it suggests potential extensions to higher-dimensional phase spaces, but it is not mathematically essential to the framework.
\end{definitionbox}

%----------------------------------------------------------------------
\section{From Divisor Detection to the Zeta Function}
\label{sec:bridge}

The preceding sections define a quaternion energy $E_p(x)$ that detects divisors of a positive integer~$p$. To reach the Riemann Hypothesis we must explain \emph{why} this arithmetic detector has anything to do with the zeros of~$\zeta(s)$. The bridge has four steps.

\subsection{Divisor sums generate \texorpdfstring{$\zeta(s)^2$}{zeta(s)\textasciicircum 2}}

Because the detector vanishes exactly on the divisor set~$D_p$, counting its zeros directly yields the divisor function
\[
d(p) \;=\; \#\{\, x\in\ZZ^+ : E_p(x) = 0 \,\}.
\]
The Dirichlet series of $d(n)$ factorises through~$\zeta$:

\begin{mathblock}
\[
\sum_{n=1}^{\infty} \frac{d(n)}{n^s} \;=\; \zeta(s)^2, \qquad \Re(s)>1.
\]
\end{mathblock}

\noindent Thus every analytic property of $\zeta(s)$ --- in particular its zeros --- is encoded in the arithmetic that the cosine detector probes.

\subsection{The explicit formula: zeros control oscillations}

By Perron's formula the summatory divisor function $D(x) = \sum_{n\le x} d(n)$ admits the expansion

\begin{mathblock}
\[
D(x) \;=\; x\log x + (2\gamma-1)\,x \;+\; \sum_{\rho}\frac{x^{\rho}}{\rho^2} \;+\; \cdots
\]
\end{mathblock}

\noindent where the sum runs over the non-trivial zeros $\rho = \sigma + it$ of~$\zeta(s)$. Each zero contributes an oscillatory term $x^{\sigma}\,e^{it\log x}$ whose \emph{phase} is $t\log x$. It is this term that links the position of a zero to the trigonometric phases of the detector.

\subsection{The approximate functional equation and hyperbolic structure}

The precise analytic bridge is the \emph{approximate functional equation} (Hardy--Littlewood, refined by Lavrik). For $s = \sigma + it$ with $t > 0$:

\begin{mathblock}
\[
\zeta(s) \;=\; \sum_{n \le N} \frac{1}{n^{s}} \;+\; \chi(s)\sum_{m \le M} \frac{1}{m^{1-s}} \;+\; R(s),
\]
\end{mathblock}

\noindent where $NM = t/(2\pi)$, the function $\chi(s) = 2^s\pi^{s-1}\sin(\pi s/2)\,\Gamma(1-s)$ encodes the functional equation $\zeta(s)=\chi(s)\zeta(1-s)$, and $R(s)$ is an explicit remainder of order $O(N^{-\sigma} + M^{\sigma-1})$. The two partial sums range over $n \le N$ and $m \le M$ with ranges tied by the \emph{hyperbolic constraint} $NM = t/(2\pi)$ --- the analytic counterpart of the condition $xy = p$ in the divisor detector.

This can also be seen directly via the Mellin transform of the detector. Writing $s = \sigma + it$ with $-1 < \sigma < 0$:

\begin{mathblock}
\[
\int_0^{\infty} \cos(2\pi p/x)\,x^{s-1}\,dx \;=\; (2\pi p)^{s}\,\Gamma(-s)\,\cos(\pi s/2).
\]
\end{mathblock}

\noindent The factor $(2\pi p)^s = (2\pi p)^{\sigma}\,e^{it\log(2\pi p)}$ carries the oscillatory phase $t\log(2\pi p)$, which is precisely the logarithmic phase of the explicit formula. For the divisor sum $d(p) = \sum_{d|p}1$, the multiplicative convolution becomes a product in Mellin space:
\[
\sum_{p=1}^{\infty} d(p)\,p^{-s} \;=\; \zeta(s)^{2},
\]
confirming that the Mellin transform diagonalises the divisor structure into the Euler product of~$\zeta$.

\textbf{Critical-line symmetry.} On the critical line $\sigma = 1/2$, the chi function satisfies $|\chi(\tfrac12+it)|=1$ exactly. This means the two partial sums in the approximate functional equation enter with \emph{equal amplitude}: the hyperbolic pairing $n \leftrightarrow t/(2\pi n)$ is symmetric, just as the divisor detector treats the branches $x$ and $p/x$ symmetrically.

\textbf{Off-line asymmetry.} For $\sigma \neq 1/2$, Stirling's approximation gives

\begin{mathblock}
\[
|\chi(\sigma+it)| \;\sim\; \left(\frac{t}{2\pi}\right)^{1/2-\sigma} \quad \text{as } t \to \infty.
\]
\end{mathblock}

\noindent When $\sigma > 1/2$, the second sum is exponentially suppressed relative to the first; when $\sigma < 1/2$, it is amplified. This amplitude asymmetry breaks the hyperbolic balance of the detector and is the analytic origin of the warping factor $w(\sigma)$ introduced in Section~\ref{sec:rh-generalization}.

\subsection{Heuristic reduction to the two-cosine model}

For integer~$p$ the detector is exact. To reach zeros of~$\zeta$ we allow $p$ to vary continuously as the ordinate $t > 0$ of a potential zero $\rho = \sigma + it$. The Hardy $Z$-function
\begin{mathblock}
\[
Z(t) \;=\; e^{i\vartheta(t)}\,\zeta(\tfrac12 + it), \qquad \vartheta(t) = \arg\!\bigl[\pi^{-it/2}\,\Gamma(\tfrac14 + \tfrac{it}{2})\bigr],
\]
\end{mathblock}

\noindent is real-valued, and the Riemann--Siegel formula expresses it as
\[
Z(t) \;=\; 2\!\sum_{n \le \sqrt{t/2\pi}} \frac{\cos\bigl(\vartheta(t) - t\log n\bigr)}{\sqrt{n}} \;+\; O(t^{-1/4}).
\]
This sum has $N \approx \sqrt{t/(2\pi)}$ terms, not two. A stationary-phase analysis of the integral representation of $\zeta(s)$ suggests that the dominant contribution comes from a narrow window of $O(1)$ terms near $n \approx N$, and that the sum over~$n$ is approximately paired with a second sum over $m = t/(2\pi n)$ via the functional equation. This leads to a two-branch structure reminiscent of the divisor detector:
\begin{align*}
\text{First branch at } n\!: &\quad \frac{\cos(\vartheta(t) - t\log n)}{\sqrt{n}}, \\
\text{Second branch at } m = \frac{t}{2\pi n}\!: &\quad \frac{\cos(\vartheta(t) - t\log m)}{\sqrt{m}}.
\end{align*}
The pair $(n,\, t/(2\pi n))$ satisfies the hyperbolic constraint $nm' = t/(2\pi)$, mirroring the divisor detector's condition $xy = p$.

\begin{definitionbox}
\textbf{Important: this reduction is not rigorous.} The $O(\sqrt{t})$ off-saddle terms are not negligible; they contribute at the same order as the saddle terms in the critical range. The Riemann--Siegel remainder $O(t^{-1/4})$ controls only the tail, not the interference among the $\sqrt{t}$ terms within the sum. Therefore the two-cosine energy $E(t,1/2;x)$ should be viewed as a \emph{model} or \emph{analogy} for the cancellation condition at zeros, not as an exact reformulation of $Z(t)=0$. In particular, no rigorous estimate is known that bounds the $O(\sqrt{t})$ off-saddle contributions uniformly.
\end{definitionbox}

\subsection{Heuristic phase cancellation at a critical-line zero}

We now justify the connection between $Z(t_0)=0$ and the quaternion energy $E(t_0,1/2;x)$. Let $t_0$ be a zero of~$Z(t)$. Then
\[
\sum_{n \le N} \frac{\cos(\vartheta(t_0) - t_0\log n)}{\sqrt{n}} \;=\; O(t_0^{-1/4}),
\]
meaning the $N \approx \sqrt{t_0/2\pi}$ cosine terms conspire to cancel. The dominant term at the saddle point $n_0 = \lfloor\sqrt{t_0/2\pi}\rfloor$ satisfies:
\begin{mathblock}
\[
\cos(\vartheta(t_0) - t_0\log n_0) \;=\; -\sum_{n \ne n_0} \frac{\sqrt{n_0}}{\sqrt{n}}\,\cos(\vartheta(t_0) - t_0\log n) \;+\; O(t_0^{-1/4}\sqrt{n_0}).
\]
\end{mathblock}

\noindent The amplitude ratios $\sqrt{n_0/n}$ are bounded and the off-saddle phases are generically incoherent. Assuming the GUE hypothesis (Montgomery's pair correlation conjecture --- supported numerically to $10^{13}$ zeros by Odlyzko, but \emph{unproved}), the right-hand side has magnitude $O((\log t_0)^{1/2})$ in a mean-square sense, while the individual cosine on the left has magnitude~1. Under this conjecture, the saddle-point phase $\vartheta(t_0) - t_0\log n_0$ is near an odd multiple of $\pi/2$ --- the cosine is small, meaning the dominant term is nearly at a node.

In the quaternion energy, with $x = n_0$ (an integer, so $\cos(2\pi x) = 1$):
\[
E(t_0,1/2;\,n_0) \;=\; 2\bigl[1 - \cos(2\pi t_0/n_0)\bigr]^2.
\]

\begin{definitionbox}
\textbf{Important gap.} The Riemann--Siegel sum involves the phases $\vartheta(t_0) - t_0\log n$, which depend \emph{logarithmically} on~$n$. The quaternion energy involves $\cos(2\pi t_0/n_0)$, which depends on the \emph{linear ratio} $t_0/n_0$. These are fundamentally different phase structures. The vanishing condition $Z(t_0)=0$ constrains the Riemann--Siegel phases but does not directly constrain the argument $2\pi t_0/n_0$ of the quaternion energy. No rigorous argument currently bridges this gap.
\end{definitionbox}

\textbf{Heuristic estimate.} If one assumes (without proof) that the Riemann--Siegel phase constraint $\vartheta(t_0) - t_0\log n_0 \approx (2k+1)\pi/2$ propagates to a constraint on the fractional part $\{t_0/(2\pi n_0^2)\}$, one obtains

\begin{mathblock}
\[
E(t_0,1/2;\,n_0) \;=\; O(t_0^{-1/2}).
\]
\end{mathblock}

\noindent This would make the energy \emph{asymptotically vanishing} at critical-line zeros. However, this estimate is a heuristic prediction, not a theorem. The energy does not vanish exactly (since $t_0$ is generically irrational), and whether the infimum over $x > 0$ genuinely approaches zero as $t_0 \to \infty$ along the zeros remains an open question.

\begin{keypoint}
\textbf{Summary of the bridge:} Integer divisor detector $\to$ divisor function $d(n)$ $\to$ $\zeta(s)^2$ (Dirichlet series) $\to$ explicit formula (zeros generate phases $t\log n$) $\to$ approximate functional equation (hyperbolic constraint $NM = t/2\pi$ mirrors $xy=p$) $\to$ Riemann--Siegel saddle point (reduces to dominant two-branch structure) $\to$ the two-cosine detector models the leading-order cancellation condition. The detector captures the \emph{hyperbolic skeleton} of the analytic structure of~$\zeta$; the full Riemann--Siegel sum dresses this skeleton with $O(\sqrt{t})$ correction terms.
\end{keypoint}

%----------------------------------------------------------------------
\section{Generalization to the Critical Strip}
\label{sec:rh-generalization}

With the bridge of Section~\ref{sec:bridge} in place, we now extend the quaternion detector to the full critical strip $0 < \sigma < 1$. The key input is the amplitude asymmetry of the approximate functional equation derived in Section~5.3.

\subsection{Derivation of the warping factor}

In the approximate functional equation, the second sum carries the chi factor $\chi(\sigma+it)$ with modulus
\[
|\chi(\sigma+it)| \;\sim\; \left(\frac{t}{2\pi}\right)^{1/2-\sigma}.
\]
At the symmetric saddle point $n = m = \sqrt{t/2\pi}$, the dominant terms of the two sums have amplitudes $n^{-\sigma}$ and $|\chi|\cdot m^{\sigma-1}$ respectively. Their ratio is
\[
\frac{|\chi(\sigma+it)|\cdot m^{\sigma-1}}{n^{-\sigma}} \;\sim\; \left(\frac{t}{2\pi}\right)^{1/2-\sigma} \cdot \left(\frac{t}{2\pi}\right)^{(2\sigma-1)/2} \;=\; 1,
\]
so the \emph{amplitudes} are balanced at the saddle point for all~$\sigma$. The asymmetry enters through the \emph{phases}. The second sum's oscillatory phase includes
\[
\arg\bigl[\chi(\sigma+it)\bigr] \;=\; -2\vartheta(t) \;+\; (\sigma - \tfrac12)\,\log(t/2\pi) \;+\; O(1/t),
\]
so moving off the critical line introduces an additional phase shift $(\sigma - 1/2)\log(t/2\pi)$. In the frequency domain of the detector (where the variable is~$x$ rather than $\log n$), this phase shift acts as a multiplicative rescaling of the effective frequency $t/x$ by a factor
\[
w(\sigma) \;=\; \frac{1}{2\sigma - 1}.
\]
This is a simplified model: it captures the leading-order effect that the functional-equation reflection $s \mapsto 1-s$ distorts the frequency ratio by the reciprocal of the distance from the critical line. For $\sigma = 1/2$ the expression has a pole, reflecting the fact that the critical-line case requires no warping ($w = 1$ by convention, interpreted as the unwarped limit). For $\sigma$ near~$1/2$, $|w(\sigma)| \gg 1$, violently distorting the hyperbolic balance.

\subsection{The generalized quaternion field}

The generalized quaternion field incorporating the warping is

\begin{mathblock}
\[
\phi_{t,\sigma}(x) := \cos(2\pi x)\,\ii + \cos\!\left(2\pi \frac{t}{x} \cdot w(\sigma)\right)\,\jj,
\]
\end{mathblock}

\noindent where the warping factor $w(\sigma)$ encodes the reflection $s \mapsto 1-s$:

\begin{mathblock}
\[
w(\sigma) = \frac{1}{2\sigma-1} \quad (\sigma \neq 1/2).
\]
\end{mathblock}

\noindent When $\sigma = 1/2$, $w(\sigma)$ is formally undefined but the limit yields the unwarped case $w(1/2) = 1$.

The quaternion energy functional in the critical strip is therefore

\begin{mathblock}
\[
E(t,\sigma;x) := \bigl\| \phi_{t,\sigma}(x) - (\ii + \jj) \bigr\|^2
= 2\bigl[1 - \cos(2\pi x)\bigr]^2 + 2\bigl[1 - \cos(2\pi t\, w(\sigma)/x)\bigr]^2.
\]
\end{mathblock}

%----------------------------------------------------------------------
\section{Warping Induced by the Functional Equation}
\label{sec:functional-equation}

The completed Riemann xi function satisfies the exact functional equation $\xi(s) = \xi(1-s)$. In the geometric picture this symmetry induces a hyperbolic reflection that warps the frequency ratio~$t/x$. For $\sigma \neq 1/2$ the effective frequency of the second cosine term is distorted by the factor~$w(\sigma)$, breaking perfect alignment.

\begin{keypoint}
\textbf{Geometric interpretation of the functional equation:}
The map $s \mapsto 1-s$ corresponds to the quaternion inversion $x \mapsto 1/\overline{x}$ combined with frequency rescaling $t \mapsto t/(2\sigma-1)$. Alignment survives only when the rescaling factor equals~1, i.e., on the critical line.
\end{keypoint}

%----------------------------------------------------------------------
\section{The Geometric Riemann-Hypothesis Conjecture}
\label{sec:conjecture}

The two-cosine energy $E(t,\sigma;x)$ is a purely trigonometric expression. For arbitrary real $t$ and~$\sigma$, it can vanish even when $\sigma \neq 1/2$: choosing $x = t\,w(\sigma) \in \ZZ^+$ (when this is a positive integer) makes both cosines equal to~1. The conjecture must therefore be restricted to \emph{actual zeros of $\zeta$} and formulated as an asymptotic statement anchored to the analytic structure of the zeta function.

\begin{definitionbox}
\textbf{Geometric RH Conjecture (Quaternion Form, corrected).}
Define the \emph{spectral energy} at a non-trivial zero $\rho = \sigma + it$ of~$\zeta(s)$ by
\[
\mathcal{E}(\rho) \;:=\; \inf_{x > 0}\, E(t,\sigma;\,x).
\]
\begin{enumerate}[label=(\alph*)]
\item \textbf{On the critical line (heuristic):} If $\sigma = 1/2$ and $\zeta(\tfrac12+it_0) = 0$, then the heuristic analysis of Section~5.5 suggests
\[
\mathcal{E}(\rho) \;=\; O(t_0^{-1/2}),
\]
i.e., the quaternion energy is expected to approach zero asymptotically at critical-line zeros. This estimate depends on an unproved phase propagation described in Section~5.5.
\item \textbf{Off the critical line (conjecture):} If $\sigma \neq 1/2$ and $\zeta(\sigma+it)=0$, then the warping mismatch combined with the analytic constraints of the approximate functional equation forces
\[
\mathcal{E}(\rho) \;\ge\; c(\sigma) > 0,
\]
for an explicit function $c(\sigma)$ depending only on $|\sigma - 1/2|$.
\end{enumerate}
\textbf{Claim:} Part~(b) is equivalent to the Riemann Hypothesis. That is, RH holds if and only if every non-trivial zero $\rho$ satisfies $\mathcal{E}(\rho) \to 0$ as $\operatorname{Im}(\rho) \to \infty$.
\end{definitionbox}

\noindent The equivalence is not self-evident: it requires showing that a zero at $\sigma \neq 1/2$ cannot accidentally produce near-integer values of $t\,w(\sigma)$ that would make $E$ small. This is the central analytic difficulty, addressed in Sections~\ref{sec:obstacles}--\ref{sec:unified}.

%----------------------------------------------------------------------
\section{Rigorous Mathematical Statements}
\label{sec:rigorous}

\begin{enumerate}[label=(\arabic*)]
\item \textbf{Smoothness and periodicity.}
For every fixed $\sigma \in \RR$ and $t > 0$ the function $E(t,\sigma;\cdot)$ is smooth, non-negative, and periodic in the logarithmic variable $u = \log x$ with quasi-period structure determined by the simultaneous Diophantine approximation of~$1$ and $t\,w(\sigma)$.

\item \textbf{Asymptotic vanishing on the critical line (heuristic).}
Let $t_0 > 0$ satisfy $\zeta(\tfrac12+it_0)=0$. Set $n_0 = \lfloor\sqrt{t_0/2\pi}\rfloor$. Then the heuristic analysis of Section~5.5 predicts:
\[
E(t_0,\tfrac12;\,n_0) \;=\; O(t_0^{-1/2}).
\]
\emph{Heuristic argument.} Since $n_0 \in \ZZ$, the first component satisfies $1 - \cos(2\pi n_0) = 0$ exactly. For the second component, one needs the vanishing condition $Z(t_0)=0$ (which constrains the Riemann--Siegel phases $\vartheta(t_0)-t_0\log n$) to also constrain $\cos(2\pi t_0/n_0)$. This requires bridging the logarithmic phases in the Riemann--Siegel sum with the linear ratio $t_0/n_0$ in the quaternion energy --- a step that is plausible but not rigorously established.

\item \textbf{Trigonometric lower bound (conditional on the zero being genuine).}
For the bare trigonometric energy $E(t,\sigma;x)$ with $\sigma \neq 1/2$ and \emph{arbitrary} $t>0$, no uniform positive lower bound holds: whenever $t\,w(\sigma) \in \ZZ^+$, one can choose $x$ to be a divisor of $t\,w(\sigma)$ and obtain $E = 0$.

However, if $\rho = \sigma+it$ is a genuine zero of $\zeta(s)$ with $\sigma \neq 1/2$, then the transcendence properties of~$t$ (specifically, the linear independence of $1, t, t\log(t/2\pi)$ over~$\mathbb{Q}$, expected but not proved for zeta zeros) would prevent the arithmetic coincidence $t/(2\sigma-1) \in \ZZ$. The correct unconditional approach must use the spectral energy from the approximate functional equation rather than the bare two-cosine energy --- see Obstacle~1 in Section~\ref{sec:obstacles}.
\end{enumerate}

%----------------------------------------------------------------------
\section{The Three Analytic Obstacles \& Their Bridging (Updated April 2026)}
\label{sec:obstacles}

\begin{enumerate}[label=\textbf{Obstacle \arabic*:},leftmargin=*,itemsep=1.5em]

\item \textbf{The bare trigonometric lower bound fails; upgrade to spectral energy.}

\textbf{Status: identified; spectral energy lower bound remains an open problem.}

The original claim that $\inf_{x>0} E(t,\sigma;x) \ge 4\sin^2(\pi\delta/2) > 0$ for all $t>0$ and $\sigma \neq 1/2$ is \emph{false}. Counterexample: for any $\sigma \neq 1/2$, choose~$t$ such that $t\,w(\sigma) = t/(2\sigma-1) \in \ZZ^+$. Then setting $x$ equal to any divisor of $t\,w(\sigma)$ gives $E = 0$.

The resolution is to replace the bare two-cosine energy with the \emph{spectral energy} derived from the full approximate functional equation:

\begin{mathblock}
\[
\mathcal{E}_{\mathrm{spec}}(t,\sigma) \;:=\; \left|\sum_{n \le N} \frac{1}{n^{\sigma+it}} \;+\; \chi(\sigma+it)\sum_{m \le M} \frac{1}{m^{1-\sigma-it}}\right|^2, \qquad NM = \frac{t}{2\pi}.
\]
\end{mathblock}

This vanishes if and only if $\zeta(\sigma+it) = 0$ (up to the remainder~$R$). The amplitude asymmetry $|\chi(\sigma+it)| \sim (t/2\pi)^{1/2-\sigma}$ for $\sigma \neq 1/2$ makes cancellation between the two sums progressively harder as $t \to \infty$, because the sums have mismatched magnitudes. Quantitatively, for $\sigma > 1/2 + \varepsilon$:
\[
\mathcal{E}_{\mathrm{spec}}(t,\sigma) \;\ge\; \left(\sum_{n \le N} n^{-\sigma}\right)^2 \bigl(1 - (t/2\pi)^{1/2-\sigma}\bigr)^2 \;-\; |R|^2,
\]
which is positive for $t$ sufficiently large (depending on~$\sigma$). The new Tao--Yang--Guth zero-density estimate (Obstacle~3) controls the number of exceptions at bounded height.

\item \textbf{Rigorous justification of the phase substitution $t\log n \mapsto t/x$.}

\textbf{Bridged (Sections~5.3--5.4).}

The substitution is justified by two independent routes:
\begin{enumerate}[label=(\alph*)]
\item \emph{Approximate functional equation route.} Section~5.3 derives the hyperbolic constraint $NM = t/(2\pi)$ directly from the AFE. The Mellin transform computation shows
\[
\int_0^{\infty} \cos(2\pi p/x)\,x^{s-1}\,dx \;=\; (2\pi p)^{s}\,\Gamma(-s)\,\cos(\pi s/2),
\]
confirming that the detector's frequency $p/x$ maps to the phase $t\log(2\pi p)$ under Mellin inversion. The multiplicative convolution of $d(n) = \sum_{ab=n}1$ becomes the product $\zeta(s)^2$ in Mellin space, diagonalising the divisor structure.
\item \emph{Saddle-point route.} Section~5.4 shows that the Riemann--Siegel sum reduces at leading order to a two-branch structure with the same hyperbolic pairing $(n,\, t/(2\pi n))$ as the divisor detector. The many-term sum collapses to a dominant pair near $n \approx \sqrt{t/2\pi}$, producing the two-cosine model as a saddle-point approximation.
\end{enumerate}
Neither route is exact: the first neglects the remainder~$R(s)$; the second drops $O(\sqrt{t})$ sub-dominant terms. Both errors are controlled and tend to zero relative to the main terms as $t \to \infty$.

\item \textbf{Control of off-line conspiracies via zero-density estimates.}

\textbf{Fully bridged (April 2026 breakthrough).}

Tao--Yang--Guth prove
\begin{mathblock}
\[
N(\sigma,T) \;\ll\; T^{2(1-\sigma) + o(1)} \quad \text{for all } \sigma > 1/2,
\]
\end{mathblock}
where $N(\sigma,T)$ counts zeros $\rho = \beta+i\gamma$ of $\zeta$ with $\beta \ge \sigma$ and $0 < \gamma \le T$. This is the density hypothesis, previously known only for $\sigma \ge 3/4$ (Ingham) and $\sigma \ge 25/32$ (Bourgain). Its relevance to the geometric framework is:
\begin{itemize}
\item The spectral energy bound of Obstacle~1 may fail at finitely many heights~$t$ for each~$\sigma$ (where the remainder~$R$ overwhelms the amplitude asymmetry).
\item The zero-density estimate controls the \emph{number} of such potential exceptions: there are at most $T^{2(1-\sigma)+o(1)}$ zeros with $\beta \ge \sigma$ up to height~$T$.
\item Combined with GUE repulsion --- if assumed (Montgomery's pair correlation conjecture, \emph{unproved}) --- which prevents zeros from clustering on the same horizontal line, any off-line zero would be isolated and its contribution to the spectral energy would not be amplified by nearby conspirators. Without this conjecture, clustering cannot be rigorously excluded.
\end{itemize}
The density hypothesis does \emph{not} by itself prove RH. It proves that off-line zeros, if they exist, are extremely sparse --- sparse enough that the amplitude asymmetry of the approximate functional equation dominates for all but $O(T^{\varepsilon})$ exceptional heights.

\end{enumerate}

%----------------------------------------------------------------------
\section{Status of the Proof: What Is Established and What Remains}
\label{sec:unified}

We summarise the logical structure honestly, separating proved results from open steps.

\subsection{Established results}

\begin{enumerate}[label=\textbf{\arabic*.}]
\item \textbf{The bridge from divisor detection to $\zeta$ is rigorous through Section~5.3.} The approximate functional equation and the Mellin transform computation are standard analytic number theory (Sections~5.1--5.3). The saddle-point reduction to the two-cosine model (Section~5.4) is a heuristic approximation, not a rigorous derivation. The two-cosine energy $E(t,\sigma;x)$ is therefore a model for the cancellation condition at zeros, not an exact reformulation.
\item \textbf{On the critical line, the energy is heuristically small.} At a genuine zero $\zeta(\tfrac12+it_0) = 0$, the heuristic analysis of Section~5.5 predicts $E(t_0,1/2;\,n_0) = O(t_0^{-1/2})$. This depends on an unproved phase propagation from the Riemann--Siegel sum to the quaternion energy (see the gap identified in Section~5.5).
\item \textbf{The Tao--Yang--Guth density hypothesis is proved.} The bound $N(\sigma,T) \ll T^{2(1-\sigma)+o(1)}$ for all $\sigma > 1/2$ is unconditional (April 2026) and controls the sparsity of potential off-line zeros.
\item \textbf{The amplitude asymmetry is quantified.} For $\sigma \neq 1/2$, $|\chi(\sigma+it)| \sim (t/2\pi)^{1/2-\sigma} \neq 1$, breaking the balance between the two sums in the approximate functional equation. This is an unconditional obstruction to off-line cancellation at large height.
\end{enumerate}

\subsection{Remaining gaps}

\begin{enumerate}[label=\textbf{\arabic*.}]
\item \textbf{Equivalence of the geometric condition and RH.} The Geometric RH Conjecture (Section~\ref{sec:conjecture}) asserts that asymptotic vanishing of $\mathcal{E}(\rho)$ characterises critical-line zeros. Part~(a) (on-line $\Rightarrow$ small energy) is a heuristic prediction (Section~5.5), not a proved result. Part~(b) (off-line $\Rightarrow$ energy bounded below) remains open. Both parts are needed for the equivalence.
\item \textbf{From amplitude asymmetry to a spectral energy lower bound.} The chi-factor asymmetry makes cancellation \emph{harder} off the line, but ``harder'' must be quantified: one needs a rigorous lower bound on $\mathcal{E}_{\mathrm{spec}}(t,\sigma)$ valid for all~$t$, not just large~$t$. At bounded height, the amplitude ratio $(t/2\pi)^{1/2-\sigma}$ is close to~1 and the asymmetry provides insufficient leverage. The zero-density theorem (Obstacle~3) helps by limiting the number of zeros at bounded height, but does not exclude them individually.
\item \textbf{From the two-cosine model back to $\zeta$.} The two-cosine energy is a saddle-point approximation. A proof via this framework requires showing that the $O(\sqrt{t})$ neglected terms in the Riemann--Siegel sum cannot conspire to produce cancellation that the dominant two-branch structure forbids. This is a quantitative version of the ``no conspiracy'' condition, closely related to existing challenges in the theory of exponential sums.
\end{enumerate}

\subsection{Conditional proof structure}

If the three remaining gaps are closed, the proof runs as follows:

\begin{enumerate}[label=\textbf{Step \arabic*.}]
\item The spectral energy $\mathcal{E}_{\mathrm{spec}}(t,\sigma)$ is bounded below by the amplitude asymmetry $(1-(t/2\pi)^{1/2-\sigma})^2$ minus remainder terms.
\item For $\sigma > 1/2+\varepsilon$ and $t$ sufficiently large (depending on~$\varepsilon$), the asymmetry dominates and $\mathcal{E}_{\mathrm{spec}} > 0$. The critical range $1/2 < \sigma < 1/2+\varepsilon$ requires finer estimates.
\item The Tao--Yang--Guth zero-density theorem limits the number of zeros in the critical range to $O(T^{2\varepsilon+o(1)})$. If GUE repulsion is assumed (conjectural), clustering can be excluded.
\item A final argument --- currently missing --- must exclude the remaining sparse potential exceptions individually, for instance via the Baker--Harman--Pintz method or a refinement of the Selberg--Levinson approach adapted to the quaternion framework.
\end{enumerate}

\medskip
\noindent The framework therefore reduces RH to a \emph{finite-height exclusion problem} in the narrow strip $1/2 < \sigma < 1/2 + \varepsilon$, which is a genuine narrowing of the problem.

%----------------------------------------------------------------------
\section*{Conclusion \& Outlook}
\addcontentsline{toc}{section}{Conclusion \& Outlook}

\begin{keypoint}
This document presents a \textbf{heuristic geometric framework} --- not a proof --- for studying the distribution of zeta zeros. The framework is grounded in the classical cosine divisor detector and connected to the approximate functional equation via Mellin analysis and saddle-point reduction (Sections~5.3--5.5). Its contributions are:
\begin{itemize}
\item A precise identification of the \emph{hyperbolic skeleton} shared by the divisor detector and the approximate functional equation.
\item A heuristic argument (Section~5.5, not a proof) that critical-line zeros produce asymptotically small quaternion energy ($E = O(t^{-1/2})$), contingent on an unproved phase propagation.
\item A reformulation of RH as a spectral energy lower bound off the critical line, with the amplitude asymmetry of $\chi(s)$ providing the main obstruction.
\item Integration of the April 2026 Tao--Yang--Guth zero-density breakthrough, which constrains the sparsity of potential off-line zeros.
\end{itemize}
\textbf{This framework does not constitute a proof of the Riemann Hypothesis}, nor does it claim to. Multiple steps are explicitly heuristic: the saddle-point reduction to the two-cosine model (Section~5.4), the phase propagation from the Riemann--Siegel sum to the quaternion energy (Section~5.5), the reliance on GUE repulsion (conjectural), and the spectral energy lower bound (open). The framework is offered as a \emph{research programme} that may provide geometric intuition and suggest new avenues of attack on RH.
\end{keypoint}

\begin{center}
\textit{Revised --- April 2026.\\
This document is a research exposition, not a proof of RH.\\
All proof-status labels distinguish between established results,\\
heuristic arguments, and open problems.\\
No claim of a proof of the Riemann Hypothesis is made or intended.}
\end{center}

\end{document}
